\section{Monte Carlo Framework}
\label{sec:monte_carlo}

This section describes the Monte Carlo component of the project. The purpose of the W3 implementation is to simulate the Heston stochastic volatility model, estimate realised variance from simulated asset paths, validate the simulation against the analytical Heston variance swap strike, and price a model-implied VIX call option. Monte Carlo simulation is used because it provides a flexible numerical approximation of risk-neutral expectations and can be applied directly to path-dependent quantities such as realised variance \citep{hirsa2024}.

The Monte Carlo engine also serves as an independent benchmark for the Fourier-cosine pricing method developed in Section~\ref{sec:cos_pricing}. While the COS method is deterministic and based on the terminal variance characteristic function, Monte Carlo uses simulated paths and therefore contains sampling and discretisation error. Comparing the two methods under the same parameter set provides a useful internal validation check.

\subsection{Simulation Objective}

The Monte Carlo implementation has three main objectives. First, it generates joint paths of the asset price \(S_t\) and variance \(v_t\) under the risk-neutral Heston dynamics. Second, it computes realised variance from simulated log returns and compares the average realised variance with the analytical variance swap strike. Third, it prices a model-implied VIX call option by transforming terminal variance into a VIX proxy and averaging discounted payoffs.

The simulation is performed under the risk-neutral measure \(\mathbb{Q}\), because derivative prices are represented as discounted risk-neutral expectations. For a terminal payoff \(\Phi\), the general pricing relation is
\begin{equation}
V_0
=
e^{-rT}
\mathbb{E}^{\mathbb{Q}}[\Phi].
\end{equation}
Monte Carlo approximates this expectation by simulating \(N\) independent paths and computing the sample average:
\begin{equation}
V_0^{MC}
=
e^{-rT}
\frac{1}{N}
\sum_{j=1}^{N}
\Phi^{(j)}.
\end{equation}

\subsection{Heston Path Simulation}

The simulated model is the risk-neutral Heston system:
\begin{align}
\frac{dS_t}{S_t}
&=
r\,dt+\sqrt{v_t}\,dW_t^S,\\
dv_t
&=
\kappa(\theta-v_t)\,dt+\sigma_v\sqrt{v_t}\,dW_t^v,\\
d\langle W^S,W^v\rangle_t
&=
\rho\,dt.
\end{align}
The correlation between the Brownian motions is implemented by generating two independent standard normal variables \(Z_v\) and \(Z_\perp\), and then defining
\begin{equation}
Z_S
=
\rho Z_v+\sqrt{1-\rho^2}\,Z_\perp.
\end{equation}
This construction ensures that
\begin{equation}
\operatorname{Corr}(Z_S,Z_v)=\rho.
\end{equation}

The time interval \([0,T]\) is divided into \(m\) equal steps of length
\begin{equation}
\Delta t=\frac{T}{m}.
\end{equation}
The simulation stores both the stock paths and the variance paths. The stock paths are used to compute realised variance, while the terminal variance values are used for model-implied VIX pricing.

\subsection{Full-Truncation Euler Scheme}

The variance process contains the square-root term \(\sqrt{v_t}\). In continuous time, the CIR process is non-negative under standard conditions. However, in discrete-time numerical simulation, a simple Euler update may produce negative variance values. Since \(\sqrt{v_t}\) is not defined for negative \(v_t\), a stability correction is required.

The implementation uses a full-truncation Euler scheme. Define
\begin{equation}
v_t^+
=
\max(v_t,0).
\end{equation}
The variance update is
\begin{equation}
v_{t+\Delta t}
=
v_t
+
\kappa(\theta-v_t^+)\Delta t
+
\sigma_v\sqrt{v_t^+}\sqrt{\Delta t}\,Z_v.
\label{eq:mc_variance_update}
\end{equation}
After the update, the next variance value is floored at zero:
\begin{equation}
v_{t+\Delta t}
\leftarrow
\max(v_{t+\Delta t},0).
\end{equation}

The stock price is simulated in logarithmic form:
\begin{equation}
\log S_{t+\Delta t}
=
\log S_t
+
\left(r-\frac{1}{2}v_t^+\right)\Delta t
+
\sqrt{v_t^+}\sqrt{\Delta t}\,Z_S.
\label{eq:mc_stock_update}
\end{equation}
The log update preserves positivity of the simulated stock price and is consistent with the diffusion form of the asset dynamics.

The simulation of square-root processes requires careful treatment near the zero boundary. Recent work on Heston and related square-root processes emphasises that discretisation choices can materially affect pricing accuracy, especially for variance-sensitive products \citep{choikwok2024heston,abijaber2024squareroot}. The full-truncation scheme is therefore used as a practical balance between numerical stability and implementation simplicity.

\subsection{Realised Variance Estimation}

After simulating stock paths, realised variance is computed from discrete log returns. For a path observed at time points
\[
0=t_0<t_1<\cdots<t_m=T,
\]
the realised variance estimator is
\begin{equation}
\widehat{\mathrm{RV}}_{0,T}^{(j)}
=
\frac{1}{T}
\sum_{i=1}^{m}
\left(
\log\frac{S_{t_i}^{(j)}}{S_{t_{i-1}}^{(j)}}
\right)^2,
\label{eq:mc_realised_variance}
\end{equation}
where \(j\) indexes the simulated path.

The Monte Carlo estimate of expected realised variance is the sample mean:
\begin{equation}
\overline{\mathrm{RV}}_{0,T}^{MC}
=
\frac{1}{N}
\sum_{j=1}^{N}
\widehat{\mathrm{RV}}_{0,T}^{(j)}.
\end{equation}
As the time discretisation becomes finer, realised variance approximates average integrated variance:
\begin{equation}
\widehat{\mathrm{RV}}_{0,T}^{(j)}
\approx
\frac{1}{T}
\int_0^T v_s^{(j)}\,ds.
\end{equation}
Therefore, the Monte Carlo mean should approximate the analytical variance swap strike:
\begin{equation}
\overline{\mathrm{RV}}_{0,T}^{MC}
\approx
K_{\mathrm{var}}(0,T).
\end{equation}

This comparison is the first validation test for the simulation engine. If the Monte Carlo mean realised variance is close to the analytical strike, the simulated variance dynamics are consistent with the theoretical expectation of integrated variance.

\subsection{Variance Swap Benchmark}

The analytical Heston variance swap strike is
\begin{equation}
K_{\mathrm{var}}(0,T)
=
\theta+
(v_0-\theta)
\frac{1-e^{-\kappa T}}{\kappa T}.
\end{equation}
The Monte Carlo engine compares this value with the mean realised variance. The error can be reported as an absolute error,
\begin{equation}
\mathrm{AbsError}_{RV}
=
\left|
\overline{\mathrm{RV}}_{0,T}^{MC}
-
K_{\mathrm{var}}(0,T)
\right|,
\end{equation}
or as a relative error,
\begin{equation}
\mathrm{RelError}_{RV}
=
\frac{
\left|
\overline{\mathrm{RV}}_{0,T}^{MC}
-
K_{\mathrm{var}}(0,T)
\right|
}{
K_{\mathrm{var}}(0,T)
}.
\end{equation}

This benchmark is useful because it checks the numerical implementation against a closed-form theoretical quantity. It does not require market data and can therefore be performed before calibration.

\subsection{Model-Implied VIX Call Pricing}

The second pricing task in the Monte Carlo engine is a model-implied VIX call. The project uses the Heston-implied VIX proxy
\begin{equation}
\widetilde{\mathrm{VIX}}_T
=
100\sqrt{A_\Delta+B_\Delta v_T}.
\end{equation}
For each simulated path \(j\), terminal variance \(v_T^{(j)}\) is transformed into terminal model-implied VIX:
\begin{equation}
\widetilde{\mathrm{VIX}}_T^{(j)}
=
100\sqrt{A_\Delta+B_\Delta v_T^{(j)}}.
\end{equation}
The call payoff is
\begin{equation}
\Phi^{(j)}
=
\left(
\widetilde{\mathrm{VIX}}_T^{(j)}
-
K
\right)^+.
\end{equation}
The Monte Carlo call price is then
\begin{equation}
C_0^{MC}(K,T)
=
e^{-rT}
\frac{1}{N}
\sum_{j=1}^{N}
\left(
100\sqrt{A_\Delta+B_\Delta v_T^{(j)}}
-
K
\right)^+.
\label{eq:mc_vix_call}
\end{equation}

This pricing method is straightforward because it only requires terminal variance values. However, since the terminal variance values come from discretised Heston paths, the result may contain both sampling error and time-discretisation bias.

\subsection{Monte Carlo Confidence Interval}

The uncertainty of the Monte Carlo estimate is measured by the standard error. Let
\[
Y_j
=
e^{-rT}
\left(
100\sqrt{A_\Delta+B_\Delta v_T^{(j)}}
-
K
\right)^+
\]
be the discounted payoff on path \(j\). The Monte Carlo price is
\begin{equation}
C_0^{MC}
=
\frac{1}{N}
\sum_{j=1}^{N}Y_j.
\end{equation}
The sample standard deviation is
\begin{equation}
s_Y
=
\sqrt{
\frac{1}{N-1}
\sum_{j=1}^{N}
\left(
Y_j-C_0^{MC}
\right)^2
}.
\end{equation}
The standard error is
\begin{equation}
SE(C_0^{MC})
=
\frac{s_Y}{\sqrt{N}}.
\end{equation}
An approximate \(95\%\) confidence interval is
\begin{equation}
C_0^{MC}
\pm
1.96\,SE(C_0^{MC}).
\end{equation}

This interval helps interpret the difference between Monte Carlo and COS prices. A small difference between the two methods is meaningful only when the Monte Carlo uncertainty is sufficiently controlled.

\subsection{Convergence Behaviour}

Monte Carlo convergence is governed by the central limit theorem. The sampling error decreases at the rate
\begin{equation}
O(N^{-1/2}),
\end{equation}
where \(N\) is the number of simulated paths. This rate is independent of the dimension of the simulated process, which is one reason Monte Carlo is widely used in finance. However, the convergence is relatively slow: reducing the standard error by a factor of two requires approximately four times as many paths.

The implementation therefore studies price stability across different path counts. In addition to increasing \(N\), time-discretisation error can be reduced by increasing the number of time steps \(m\). In practice, both sources of error must be considered. A large number of paths does not eliminate bias caused by a coarse time grid.

\subsection{Interpretation of W3 Results}

The W3 component provides a flexible simulation-based pricing framework. Its first validation step is the comparison between simulated realised variance and the analytical variance swap strike. Its second output is a Monte Carlo price for a model-implied VIX call. These two tasks show that the engine can be used both for path-based quantities and terminal-payoff derivatives.

The Monte Carlo engine is not intended to be the fastest pricing method in the project. Instead, its role is to provide a transparent and flexible benchmark. It is especially useful because the simulated paths can later be reused for additional tasks such as scenario analysis, preliminary variance-risk comparisons, and empirical extensions once market data are added.

\subsection{Limitations of the Monte Carlo Engine}

The current Monte Carlo implementation has several limitations. First, it uses an Euler-type discretisation rather than exact simulation of the CIR process. This is practical and stable, but it may introduce discretisation bias. Second, the current implementation uses synthetic parameters rather than calibrated market parameters. Third, the VIX quantity priced in the simulation is a Heston-implied proxy rather than the official Cboe VIX index.

These limitations do not invalidate the W3 implementation. Rather, they define the scope of the current stage. The goal of W3 is to build and validate a simulation engine under controlled model assumptions. Future stages can extend this engine by adding market calibration, exact CIR sampling, variance risk premium backtesting, and real-data validation.